\documentclass[11pt,a4paper]{article}

% -------------------- PACKAGES --------------------
\usepackage[margin=1in]{geometry}
\usepackage{amsmath,amssymb}
\usepackage{graphicx}
\usepackage{booktabs}
\usepackage{xcolor}
\usepackage{hyperref}
\usepackage{setspace}
\usepackage{fancyhdr}
\usepackage{lastpage}
\usepackage{authblk}
\usepackage{tcolorbox}
\usepackage{tocloft}

% -------------------- PAGE STYLE --------------------
\pagestyle{fancy}
\fancyhf{}
\cfoot{\thepage\ / \pageref{LastPage}}
\renewcommand{\headrulewidth}{0pt}


% -------------------- COMMENT BOX STYLE --------------------
\newtcolorbox{reviewbox}{
  colback=green!8!white,
  colframe=green!60!black,
  boxrule=0.7pt,
  arc=2pt,
  left=8pt,
  right=8pt,
  top=6pt,
  bottom=6pt
}

\newtcolorbox{updatebox}{
  colback=blue!6!white,
  colframe=blue!60!black,
  boxrule=0.7pt,
  arc=2pt,
  left=8pt,
  right=8pt,
  top=6pt,
  bottom=6pt
}

% NEW: author-initiated change (not requested by a reviewer)
\newtcolorbox{authorbox}{
  colback=orange!8!white,
  colframe=orange!70!black,
  boxrule=0.7pt,
  arc=2pt,
  left=8pt,
  right=8pt,
  top=6pt,
  bottom=6pt
}

% NEW: open item that still needs the author's attention
\newtcolorbox{flagbox}{
  colback=red!6!white,
  colframe=red!70!black,
  boxrule=0.7pt,
  arc=2pt,
  left=8pt,
  right=8pt,
  top=6pt,
  bottom=6pt
}

% -------------------- HIGHLIGHT --------------------
\newcommand{\rev}[1]{\colorbox{yellow}{#1}}

% -------------------- TITLE --------------------
\title{2nd Round Responses to Reviewers’ Comments for Manuscript\\[0.8em]
\textbf{Implementation of Communication and Group Formation in the Agent Based Modelling of Forced Migrants: Mali Case Study}\\[0.8em]
Addressed Comments for Publication to\\
\textbf{The Journal of Artificial Societies and Social Simulation (JASSS)}
}

\author[1]{Edward King}
\author[1]{He-In Cheong}
\author[1]{Arnab Majumdar}
\affil[1]{Department of Civil and Environmental Engineering, Imperial College London\\
South Kensington Campus, London SW7 2AZ}

\date{}

\begin{document}

% ======================================================
% ===================== TITLE PAGE =====================
% ======================================================

\maketitle
\vfill

\begin{center}
\textbf{Response Letter}
\end{center}

\newpage

% ======================================================
% ===================== CONTENTS =======================
% ======================================================

\setcounter{page}{1}
\renewcommand{\contentsname}{Contents}
\tableofcontents

\newpage

% ======================================================
% ===================== PREFACE ========================
% ======================================================

\section*{Preface}
\addcontentsline{toc}{section}{Preface}

We thank the editor and both reviewers for their careful second-round comments. Every comment has been
addressed below, and where two reviewers raised the same point (e.g.\ the parenthetical ``unethical'' in
Section~1.3, the ``explain'' typo, and the framing of Section~1.9) the changes have been \emph{merged} so
that a single, consistent revision appears in the manuscript.

Alongside the review responses, we have made a set of substantive, author-initiated revisions that improve
the rigour, reproducibility and robustness of the study. The most important of these is that
\textbf{all simulations have been re-run with a fixed random seed}, so that the results reported in the
manuscript are fully deterministic and reproducible rather than subject to run-to-run stochastic variation.
As a consequence, \textbf{every numerical result, table and figure has been regenerated}, and the reported
figures now differ from the previous submission (for example, headline camp precision is now 77.85\%
rather than 75.1\%). We have also restricted all reported accuracy, precision and relative-difference
metrics to \textbf{empirical UNHCR observation dates only} (interpolated values are used for visualisation
but never for scoring), which is a more honest basis for evaluation. Additional engineering improvements
support this: a fix enabling the model to run on Windows, updated dependencies (e.g.\ \texttt{pandas}),
alternating-slot checkpointing so long runs can be resumed after an interruption, and a simplified mechanism
for toggling the grouping, communication and roulette modules on and off (which underpins the ablation study).

To keep this document transparent, we distinguish four kinds of content by colour:

\vspace{0.5em}
\noindent\textbf{Note:} The light \textbf{green} boxes reproduce the original reviewer comments. The
\textbf{blue} boxes give the corresponding revised text as it now appears in the manuscript. The
\textbf{orange} boxes document author-initiated changes that were \emph{not} requested by a reviewer, together
with the reason for the change (Section~\ref{sec:author}). The \textbf{red} boxes flag items that still
require an author decision or input before final submission (collated in Section~\ref{sec:open}).

\vspace{0.5em}
\noindent Comments on the appendices are not itemised here; the appendix figures were regenerated to match
the reseeded run, but no other appendix changes are documented.

\newpage

% ====================================================
\section*{Response to Reviewer 1}
% ====================================================
\addcontentsline{toc}{section}{Response to Reviewer 1}

\subsection*{Abstract}

\begin{reviewbox}
Remove the explicit mention of the algorithm for network distances, it is not innovative enough that you would like to mention it at the beginning of the paper. Just say this, without the brackets:
\end{reviewbox}

\textbf{Response:} Agreed. The parenthetical ``(via Dijkstra's algorithm)'' has been removed.

\vspace{0.3cm}
\begin{updatebox}
The model incorporates both physical and network-based path distances between nodes and uniquely
integrates dynamic border status changes over time, capturing real-world movement constraints and
resistance across international borders.
\end{updatebox}

\vspace{0.3cm}
\begin{flagbox}
\textbf{Note (author-initiated, cross-reference):} the abstract also now reports \textbf{77.85\%} average
daily precision (was 75.1\%), grouping improving precision by \textbf{48.8\%} (was 46\%), and communication
reducing relative difference by \textbf{61.26\%} (was 56\%). These are consequences of the fixed-seed re-run
and the empirical-date-only metric convention; see Section~\ref{sec:author}, item~A1--A2.
\end{flagbox}

\subsection*{Introduction}

\subsubsection*{Section 1.3}

\begin{reviewbox}
Drop the word ``unethical'', I think what you wanted to say is that ``engineering population changes by using forced migration would constitute unethical behaviour/policies''; however, that section does not claim that either ``you, the researchers'' or ``governments'' do it, but simply that forced migration affects demography. This argument is uncontroversial, and does not have an ethical dimension.
\end{reviewbox}

\textbf{Response:} That was correct in what we were trying to say. However, as you pointed out it is
redundant in this paragraph (and research). It has thus been removed. \emph{(Reviewer~2 raised the same
point; the two comments are handled with this single change.)}

\vspace{0.3cm}
\begin{updatebox}
The dynamics of forced migration are uncertain and highly complex, they are considered the most un-predictable
drivers of population change \citep{NRC2000}.
\end{updatebox}

\subsubsection*{Section 1.9}

\begin{reviewbox}
Maybe replace ``communication networks'' with ``social networks'' because it is ambiguous, it may indicate ``telecommunications'' the way it is expressed.
Also ``expalin'' to ``explain''.
\end{reviewbox}

\textbf{Response:} We agree with the ambiguity. ``Communication networks'' has been replaced with
``social networks'' and the spelling error corrected. \emph{This paragraph now also incorporates
Reviewer~2's request to specify what ``improved predictive outcomes'' refers to (see Reviewer~2, Section~1.9);
both edits are merged into the single sentence below, so the final paragraph reads:}

\vspace{0.3cm}
\begin{updatebox}
The aim of this research is to enhance the accuracy of an ABM for forecasting the movement of forced migrants
by integrating group dynamics and \textbf{social networks}. Using the case study of Mali in 2012, the research
seeks to demonstrate that a more granular understanding of the forced migrant experience leads to improved
predictive outcomes, particularly in modelling camp population distributions, destination selection, and the
temporal dynamics of displacement, through the incorporation of grouping mechanisms and communication
processes, compared to baseline models such as Flee.

\medbreak
\textbf{Also in Appendix A:} The range of interaction is therefore defined both by explicit social networks
and by emergent group membership. These assumptions reflect empirical observations that displaced populations
rely heavily on word-of-mouth and kinship ties to navigate uncertain environments.

\medbreak
For brevity, the Supplementary Notes will be used to explain concepts and assumptions referenced in the main
text in higher detail. \emph{(``expalin'' $\rightarrow$ ``explain'' corrected; Reviewer~2 flagged the same typo.)}
\end{updatebox}

\subsection*{Literature Review}

\subsubsection*{Section 2.1}

\begin{reviewbox}
Replace this:

``Its validation focuses on final outcomes, overlooking key decision-making processes, which limits its ability to capture the full complexity of forced migrant movements (Suleimenova et al. 2017).''

With this:

``The model, insofar as it does not consider the interaction between agents and groups, is missing a key component of the real-world behaviour of agents, whose choices are affected by the choices of other agents''.
\end{reviewbox}

\textbf{Response:} We thank the reviewer for the clearer phrasing, which resolves the ambiguity of the
previous wording.

\vspace{0.3cm}
\begin{updatebox}
The model, insofar as it does not consider the interaction between agents and groups, is missing a key
component of the real-world behaviour of agents, whose choices are affected by the choices of other agents.
\end{updatebox}

\subsection*{Methodology}

\subsubsection*{Section 3.3}

\begin{reviewbox}
You wrote a footnote that defines what is an IDP and what is a refugee. Why don't you use it in section 3.3.? When you talk about ``refugees hosted in Mali'', do you mean Malian citizens or foreigners? If the former, replace ``refugees'' with ``IDPs''.
\end{reviewbox}

\textbf{Response:} The distinction is now applied consistently in Section~3.3, and ``refugees hosted in Mali''
is clarified to mean \emph{foreign} nationals, in line with the footnote definitions.

\vspace{0.3cm}
\begin{updatebox}
At the heart of the Central Sahel, Mali faces multifaceted challenges including inadequate security, conflict, and
climatic shocks. According to \citet{unhcr_refugee_statistics_2025} statistics, in 2012 (the year of the case study)
Mali hosted approximately 14,000 foreign refugees while about 150,000 Malian nationals were displaced abroad as refugees
\citep{unhcr_mali}. By 2024, these figures had risen to around 136,000 foreign refugees hosted in Mali and 354,000 Malian
refugees residing in other countries \citep{unhcr_mali}. Beyond displacement, Mali grapples with severe environmental
hazards such as rising temperatures, projected to increase by 2.5\textdegree C by 2080 \citep{76-wachiaya2021warming}, as well as
risks of deforestation, desertification, and soil erosion that compromise water availability. Human trafficking also
remains prevalent along migration routes, with digital technologies playing a dual role as both a risk factor and a
resource for safety information \citep{UNHCR2022Trafficking}. Regionally, ongoing conflicts in the Central Sahel have
led to more than 410,000 refugees and approximately 2.1 million internally displaced persons \citep{80-cheshirkov2022decade}.
\end{updatebox}

\subsubsection*{Section 3.4}

\begin{reviewbox}
When you talk about the variable ``wealth distribution'' it is not clear what you mean and how you built it. \ldots{} Please expand and clarify how exactly this variable is built, it remains unclear.
\end{reviewbox}

\textbf{Response:} We have separated the description of the data sources from the model implementation.
Section~3.4 now only describes the datasets used (including the World Inequality Dataset), while the full
functional form of the wealth variable and its use inside the ABM are given explicitly in Section~4.5. This
improves clarity and reproducibility.

\vspace{0.3cm}
\begin{updatebox}
\textbf{Section 6.1 (data-source framing):} GRASP is an ABM designed to predict immediate forced migrant
flows after sudden conflicts, utilising empirical evidence on the grouping and communication of forced
migrants. Although applied to Mali 2012 with UNHCR micro-data \citep{He-In}, GRASP can be extended to other
cases using a combination of qualitative and quantitative data, including conflict event data
\citep{243-acled2024data}, demographic structure (age and gender distributions), settlement and camp
population data, and income inequality information derived from the World Inequality Dataset.

\medbreak
\textbf{Section 4.5 (implementation):} In the model, wealth influences agents' transport speed and destination
selection. Wealth is assigned using a percentile-based transformation of the form:
\begin{equation}
W(p) = 100000 \cdot (100 - p)^{-0.25} - 30000,
\end{equation}
where \(p \in [0,100]\) denotes the agent's wealth percentile. Age affects both mobility and dependent
classification (relevant to group formation), while gender is not currently used but remains a possible
future extension.
\end{updatebox}

\vspace{0.3cm}
\begin{reviewbox}
Sudan is ok, Moldova and Myanmar are not good comparisons. Add a sentence that states something like ``For the purpose of this particular paper, which focuses more on the methodology than on the exact replication of migrant distributions for Mali, we assume that migrants in Mali behave similarly to those from Moldova or Myanmar''.
\end{reviewbox}

\textbf{Response:} Added. The behavioural assumption is now stated explicitly.

\vspace{0.3cm}
\begin{updatebox}
Micro-data were extracted from UNHCR case studies on Sudan, Moldova, and Myanmar by \citet{He-In}. Household size
distributions, though not directly related to Mali, are comparable to Sudan due to similar mean household sizes
\citep{dhs_sr261,ESRI_AverageHousehold}. Moldova and Myanmar were the only available sources for data on forced migrant
travel groups and their sources of information. For the purpose of this paper, which focuses on the methodological
framework rather than the exact replication of migrant distributions in Mali, we assume that migrants in Mali behave
similarly to those observed in the Moldova and Myanmar case studies.
\end{updatebox}

\subsubsection*{Section 3.11}

\begin{reviewbox}
The confusion matrix seems superfluous, and it occupies half of a page, I would personally remove it but it's your choice. Virtually all readers of the paper will understand the terminology contained in it regardless of your picture.
\end{reviewbox}

\textbf{Response:} We agree the confusion-matrix figure can be assumed familiar to the typical reader (further
explanation also remains in the supplementary note), so the figure has been removed and the paragraph updated.

\vspace{0.3cm}
\begin{updatebox}
Estimating the number of forced migrants can be viewed as a form of binary classification (see Supplementary note 5).
For this reason, in our analysis, we will focus on precision and accuracy. Precision is the ratio of correctly
classified positive outcomes to total predicted positive outcomes, while accuracy is the total number of correctly
predicted outcomes \citep{confusion}.
\end{updatebox}

\vspace{0.2cm}
\begin{flagbox}
\textbf{Knock-on effect:} removing this figure renumbered every subsequent figure (v1 Fig.~2 confusion matrix
is gone), so v1 Fig.~3\ldots{}14 are now Fig.~2\ldots{}13. Please verify all \texttt{\textbackslash ref} figure
cross-references resolve after recompilation.
\end{flagbox}

\subsubsection*{Section 3.13 and Section 3.14}

\begin{reviewbox}
Sections 3.13 and 3.14 seem unnecessary. You can simply indicate the measure for accuracy that you chose.
\end{reviewbox}

\textbf{Response:} Partially actioned, but \textbf{left open per author decision.} In the current manuscript the
two v1 paragraphs (v1~\S3.13, the ``precision of 100\%'' passage, and v1~\S3.14, the camp-proportion
normalisation passage) have in fact been \emph{removed}, and their footnotes with them. New methodological
paragraphs were added in their place (see Section~\ref{sec:author}, item~A2). What is \emph{not} yet decided is
whether to retain a short sentence justifying why precision matters more than accuracy in this application.

\vspace{0.3cm}
\begin{flagbox}
\textbf{OPEN --- decision needed.} Do you want to (a) leave §3.13/§3.14 fully removed, or (b) add one short
sentence explaining why precision is prioritised over accuracy here (many true negatives inflate accuracy)?
You noted you were unsure what the reviewer intended, so this is left for you. If (b), the old footnote
(``Averaged accuracy is not relevant due to many true negatives skewing the data'') is a ready one-line
justification.
\end{flagbox}

\subsection*{Agent Based Models}

\subsubsection*{Section 4.9}

\begin{reviewbox}
Agents who are just born are unlikely to be able to travel alone \ldots{} so births in your model essentially indicate ``an agent has come to age and can in principle move autonomously, if selected to do so; otherwise, they travel with their family''.
\end{reviewbox}

\textbf{Response:} The treatment of newborns is now stated explicitly: they are assigned to a family group and
make no independent migration decisions, and the absence of an age-based transition to autonomy is noted as a
current limitation.

\vspace{0.3cm}
\begin{updatebox}
Family sizes are defined according to micro-data from Sudan\footnote{Ref}. Families are formed at the start of the
simulation, with the possibility of increasing in size through births and decreasing through deaths. Newborn agents
are immediately assigned to a family group and do not make independent migration decisions; instead, their movement
is fully determined by the behaviour of the group to which they belong. While births are included primarily to support
long-term simulations and group cohesion, their impact is minimal over shorter time frames. A limitation of the current
model is that there is no age-based transition to independent decision-making. Deaths, however, play a more active
role---especially in conflict zones---where they influence both agent survival and the dissemination of risk
information to other agents.
\end{updatebox}

\vspace{0.2cm}
\begin{flagbox}
\textbf{OPEN --- placeholder reference.} The Sudan footnote is still the literal placeholder ``Ref''
(\texttt{\textbackslash footnote\{Ref\}}). Please insert the real citation (this is present in both the paper
and this letter).
\end{flagbox}

\subsubsection*{Section 4.20}

\begin{reviewbox}
In section 4.20, you could consider adding a point at the end that states ``in principle, if real-world trajectories sampled at daily frequency were available, these beta parameters could be learnt''. \ldots{} [reinforcement learning suggestion for a future paper].
\end{reviewbox}

\textbf{Response:} Added. We note the possibility of learning the $\beta$ parameters from trajectory data, and
expand on reinforcement-learning / gradient-based calibration in the Future Research section, while leaving a
full implementation to future work.

\vspace{0.3cm}
\begin{updatebox}
In principle, if real-world trajectories sampled at a daily frequency were available, these $\beta$ parameters
could be learned directly from data.

\medbreak
\textbf{Added to Future Research:} Future developments in computational intelligence methods, such as data
mining and reinforcement learning, could reflect realistic agent behaviour, provided they are justified and
not driven by trends. \textbf{In particular, reinforcement learning or gradient-based approaches could be
explored to calibrate behavioural parameters (e.g.\ $\beta$ weights) from observed or synthetic trajectory
data.} The goal is to balance the improvement in granularity with the increase in computational complexity.
\end{updatebox}

\vspace{0.2cm}
\begin{authorbox}
\textbf{Related author-initiated change:} the utility equation (now Eq.~12) additionally gained a bias term
$\beta_0$, with $\beta_0=150$ added to Table~2 and $\beta_{F,i}$ changed from $1/2$ to $1$. See
Section~\ref{sec:author}, item~B1.
\end{authorbox}

\subsection*{Results}

\subsubsection*{Figure 6}

\begin{reviewbox}
Here are the images obtained, which do not match those presented in the paper:
\href{https://pasteboard.co/LzLPAUiGF1m9.png}{Image 1} \href{https://pasteboard.co/KUbKqpYE6Ia1.png}{Image 2}.
Additionally, the file \texttt{Country\_split\_refugees\_w0\_roulette.csv} is not present in the repository. \ldots{} the repository should allow an interested reader to clone it, run all scripts, and reproduce the figures exactly as presented in the paper.
\end{reviewbox}

\textbf{Response (draft --- please verify before submission):} We thank the reviewer for testing reproducibility.
The results have now been fully regenerated using a \emph{fixed random seed}, so a clean run reproduces the
reported figures deterministically rather than sampling a new stochastic realisation each time. Alongside this,
the analysis scripts and the output files required to regenerate every figure have been added to the repository.

\vspace{0.3cm}
\begin{flagbox}
\textbf{OPEN --- must confirm before final submission.} (1)~Confirm that
\texttt{Country\_split\_refugees\_w0\_roulette.csv} (and any other missing outputs) is now committed to the
repository. (2)~Do a fresh clone on a clean machine and confirm all scripts run end-to-end and reproduce the
figures \emph{as printed}. (3)~State the exact random seed and how to set it in the README. Until (1)--(3) are
verified I cannot assert reproducibility on your behalf.
\end{flagbox}

\subsubsection*{Section 5.5}

\begin{reviewbox}
Section 5.5. has a problem, fix it before final release, let us pretend nobody noticed it.
\end{reviewbox}

\textbf{Response:} The stray editor's note (``Change this and all other section references to be links'') has been
removed. The related placeholder near Section~6.9 has also been addressed, and proper cross-reference links added
in both places. The paragraph now records that final results use a fixed random seed for reproducibility.

\vspace{0.3cm}
\begin{flagbox}
\textbf{OPEN --- typo in fix.} The replacement text points to ``Section~\textbf{66}.12'' (doubled digit).
Correct to Section~6.12 (ideally a \texttt{\textbackslash ref}).
\end{flagbox}

\subsubsection*{Section 5.13}

\begin{reviewbox}
section 5.13 is too long, please cut it in half. I would remove everything after footnote 12.
\end{reviewbox}

\textbf{Response:} Agreed --- little analysis was added beyond footnote~12, so that portion has been removed.
The paragraph was also updated to the re-run figures (it now discusses Tabareybarey and Mbera rather than
Mentao and Abala; see Section~\ref{sec:author}, item~C1).

\vspace{0.3cm}
\begin{updatebox}
Fig.~\ref{flee0} compares Flee and GRASP for Tabareybarey and Mbera (with all 7 comparable camps in Fig.~S7 of the
Supplementary Information). When scaled in the same fashion as Flee, it becomes apparent that GRASP is unable to
predict forced migrant flow with a higher degree of precision than Flee when assessed solely by relative
differences. However, GRASP more consistently captures the characteristic temporal behaviour of refugee movements,
identifying the timing of major influxes, periods of population stability, and the onset of population decline while
providing a reasonable estimate of camp populations. Large gradients are also detected, which are crucial for
humanitarian support as they indicate when immediate aid is required for a camp. Regarding the characteristic shape
of the forced migrant flows, Tabareybarey, Mbera, Abala, Mangaize, and Mentao all exhibit population trends that more
closely resemble the empirical observations than those produced by Flee, even where discrepancies in population
magnitude remain\footnote{This is particularly important for humanitarian operations if refugee registrations are
incomplete or delayed.}.
\end{updatebox}

\subsubsection*{Figure 12}

\begin{reviewbox}
Ok I understand you don't want to change ylim in the top figure, but you really cannot keep half the chart empty in the bottom figure. Please fix it by rescaling it and limit to y=1.25.
\end{reviewbox}

\textbf{Response:} The relative-difference figure was regenerated with the re-run data (now the Tabareybarey/Mbera
pair).

\vspace{0.3cm}
\begin{flagbox}
\textbf{OPEN --- not yet fixed.} In the regenerated figure the lower subplot still spans $y=0$ to $2.0$ with the
upper half empty. Please cap the lower subplot at $y=1.25$ as requested and then write a one-line confirmation
here.
\end{flagbox}

\subsection*{Discussion}

\subsubsection*{Section 6.2}

\begin{reviewbox}
`` leading to unrealistic simulation'' is repeated.
\end{reviewbox}

\textbf{Response:} Corrected --- the duplicated clause has been removed. Thank you for spotting it.

\vspace{0.3cm}
\begin{updatebox}
GRASP addresses limitations in current forced migration ABMs like the Flee 1.0 model, which fail to account for forced
migrant experiences and micro-processes, leading to simulations with unverified fidelity. This research assesses the
effectiveness of incorporating migrant topology (the interactive element of agents between themselves and their
environment), focusing on grouping and communication, to improve forecasting in forced migration ABMs. These aspects
were chosen based on advancements in literature \citep{84-hinsch2019rumours,163-bijak2021routes,121-collins2016agent-based}.
The dynamic weighting of nodes, based on evolving conflict data and border accessibility, also allows the model to
reflect real-time constraints in migrant decision-making.
\end{updatebox}

\subsubsection*{Section 6.9}

\begin{reviewbox}
Fix the missing link if one is supposed to go there, otherwise remove it if it is a placeholder.
\end{reviewbox}

\textbf{Response:} Addressed together with Section~5.5 above --- the reference link has been fixed.

\vspace{0.3cm}
\begin{flagbox}
\textbf{OPEN --- broken cross-reference nearby.} Section~6.10 currently reads ``as noted in Section~5~\textbf{??}'';
the \texttt{\textbackslash ref} did not resolve. Please point it at the correct results section.
\end{flagbox}

\subsubsection*{Section 6.12}

\begin{reviewbox}
I would like to understand how exactly expensive is it \ldots{} Is it one laptop full time for one week? Or did you have to rent a cluster of GPUs for a year? You should insert in this section a clarification on the actual time it took you to run everything \ldots
\end{reviewbox}

\textbf{Response (template --- needs your numbers):} We agree a concrete figure is warranted. GRASP is CPU-bound
(no GPUs) and runs on a single machine; long runs are now resumable via checkpointing.

\vspace{0.3cm}
\begin{flagbox}
\textbf{OPEN --- insert real runtime.} Please state the actual cost of a full \texttt{frac=100} run, e.g.\
``approximately \emph{X} hours on a single laptop (\emph{n} cores, \emph{y}\,GB RAM), and roughly \emph{Z} hours
for the full ablation set.'' A one-sentence quantitative answer here will fully resolve the comment.
\end{flagbox}

\subsubsection*{Section 6.14}

\begin{reviewbox}
Add a clarification that the reduction of 100 people to 1 agent may introduce distortions, when we consider that intra-family dynamics between agents are also present.
\end{reviewbox}

\textbf{Response:} Added --- aggregating 100 individuals into a single agent may distort intra-family dynamics.

\vspace{0.3cm}
\begin{updatebox}
Additionally, this level of aggregation may introduce distortions in intra-family dynamics, as interactions that
would occur within families are instead represented between agents.
\end{updatebox}

\subsubsection*{Section 6.18}

\begin{reviewbox}
You do have some type of wave-like behaviour \ldots{} Please add a statement of the form ``however, our model can partially represent dependencies in the selection of migration trajectories by future agents on the basis of the choices by past agents insofar as the arrival at destination by one migrant can influence the future arrival to that same destination of other migrants''.
\end{reviewbox}

\textbf{Response:} Added --- the model partially captures such wave-like dependencies through destination-based
interactions.

\vspace{0.3cm}
\begin{updatebox}
However, the model can partially represent dependencies in the selection of migration trajectories, insofar as the
arrival of one agent at a destination can increase the likelihood that subsequent agents select the same destination.
\end{updatebox}

\subsubsection*{Section 6.22}

\begin{reviewbox}
Remove this: ``To improve scalability and deployment speed, future iterations of GRASP could explore unsupervised methods, such as minimum spanning trees or distance-constrained heuristics \ldots{} before manual refinement introduces higher spatial and contextual accuracy''.
\end{reviewbox}

\textbf{Response:} Removed. We agree it added little to an already lengthy Discussion and lacked grounding in the
paper's theme. (No replacement text; the surrounding paragraph now ends at the ``overly dense and unrealistic
networks'' sentence.)

\subsubsection*{Section 6.23}

\begin{reviewbox}
It would make sense to use a sentence of the form ``the agent does not have global memory of the trajectory, but only a local memory of its immediate past; a future iteration of this study may include a long-term memory for the agent that is incrementally populated as the agent migrates''. \ldots{} Also drop the sentence ``While this would increase memory and computational load'' \ldots
\end{reviewbox}

\textbf{Response:} The local/global memory distinction and the scope for a longer-term, incrementally populated
memory have been added, and the note on computational cost removed.

\vspace{0.3cm}
\begin{updatebox}
The current model includes a limited form of memory, where agents retain and share knowledge of previously encountered
conflict zones through a system of ``no-go'' nodes. This mechanism allows agents to avoid revisiting dangerous areas and
contributes to more informed route planning. The agent does not have global memory of its full trajectory, but instead
operates with a local memory of its immediate past interactions and exposures. However, this memory is specific to
conflict exposure and does not extend to a broader recollection of the agent's full path or personal experiences during
migration. A future iteration of this study may incorporate a longer-term memory that is incrementally populated as the
agent migrates. Future enhancements could build on this foundation by introducing fuller path-dependency, where agents
remember the sequence of nodes traversed or assign qualitative experiences (e.g., hardship, delay) to locations. This
would allow agents to refine decision-making based on accumulated experiences, such as avoiding inefficient or
distressing routes, even in the absence of direct conflict.
\end{updatebox}

\subsubsection*{Section 7.3}

\begin{reviewbox}
Remove the bullet list and write in a single paragraph.
\end{reviewbox}

\textbf{Response:} The bulleted list has been collapsed into a single, lightly reworded paragraph.

\vspace{0.3cm}
\begin{updatebox}
Key findings indicate that even at a 1\% population scale, GRASP outperforms the Flee model in several ways. It is able to
estimate the number of refugees in camps directly from conflict data with improved precision (a capability that Flee lacks
without relying on refugee registration data). It also identifies sudden gradients, peaks, and periods of stability in UNHCR
camp populations, and classifies not only refugees but also IDPs and returnees with greater precision. Additionally, GRASP
utilises empirical data for the micro-processes involved in refugee movement where possible, rather than relying on assumed
proxies, and provides a utility-based decision model that can be further refined through interviews with Malian residents and
refugees. Finally, it is capable of estimating the number of non-camping refugees, including those seeking refuge in
neighbouring settlements and those fleeing via airports.
\end{updatebox}

\subsection*{Bibliography}

\begin{reviewbox}
This reference is wrong, you inserted probably ``Article in journal'', while it is a ``Report'': Network, G. O. (2017). Global opportunity report 2017. Tech. rep., DNV GLAS. \ldots{} please fix it since otherwise it seems like Network is somebody's surname: https://www.dnv.com/publications/global-opportunity-report-2017-84879/
\end{reviewbox}

\textbf{Response:} Corrected in the \texttt{.bib} to a technical-report entry so the author no longer renders as
``Network''.

\vspace{0.3cm}
\begin{flagbox}
\textbf{OPEN --- verify after compile.} The citation in §1.1 still renders as ``(Network 2017)''. Recompile and
confirm it now reads as the DNV GL report (correct author/organisation and ``Tech.\ rep.'' type), then delete
this flag.
\end{flagbox}

% ====================================================
\section*{Response to Reviewer 2}
% ====================================================
\addcontentsline{toc}{section}{Response to Reviewer 2}

\subsection*{Introduction}

\subsubsection*{Section 1.3}

\begin{reviewbox}
I don't know that the parenthetical ``and unethical'' is necessary here. I had to read it three times because it made the paragraph less clear.
\end{reviewbox}

\textbf{Response:} Agreed (as did Reviewer~1). The bracketed part has been removed --- handled with the single
change shown under Reviewer~1, §1.3.

\vspace{0.3cm}
\begin{reviewbox}
I'm unclear on the evidence to claim that ``migration modelling itself remains underdeveloped'' and that ``models often struggle to produce reliable theoretical predictions.'' - these warrant citations.
\end{reviewbox}

\textbf{Response:} Rewritten to reflect the literature and supported with citations (Massey et al.\ 1993;
Davenport et al.\ 2003; Van Hear et al.\ 2018; Cheong 2023), reintroduced from an earlier draft.

\vspace{0.3cm}
\begin{updatebox}
However, several studies highlight limitations in current migration modelling approaches, including weak theoretical
foundations, methodological biases, and insufficient treatment of complex driver interactions (Massey et al.\ 1993;
Davenport et al.\ 2003; Van Hear et al.\ 2018; Cheong 2023).
\end{updatebox}

\subsubsection*{Section 1.4}

\begin{reviewbox}
Citation for ``new levels of granularity'' is from 1990. Demographics modeling has progressed in granularity since that time, I believe.
\end{reviewbox}

\textbf{Response:} Reworded to present granularity as an ongoing development, with a contemporary reference
(Bijak 2022) added.

\vspace{0.3cm}
\begin{updatebox}
With the advancement in conceptualising population growth and decline mechanisms, improvement of modelling techniques, and
rapid development of computational power, demographic forecasting has evolved to increasingly fine levels of granularity
\citep{106-willekens1990demographic, bijak2022}.
\end{updatebox}

\vspace{0.2cm}
\begin{authorbox}
\textbf{Related author-initiated change:} the sentence ``These efforts focus on bridging the gap between statistical
information and the understanding of causal systems'' was also dropped from §1.4 (tightening only). See
Section~\ref{sec:author}, item~C4.
\end{authorbox}

\subsubsection*{Section 1.5}

\begin{reviewbox}
I'm not sure I follow what this paragraph is trying to say about ``researchers in host countries.''
\end{reviewbox}

\textbf{Response:} The paragraph now closes by linking the ``methodological nationalism'' point to the modelling
limitations it introduces.

\vspace{0.3cm}
\begin{updatebox}
These limitations can constrain the development of more comprehensive and representative migration models.
\end{updatebox}

\subsubsection*{Section 1.6}

\begin{reviewbox}
Literature on push/pull factors and claim that ``theoretical understanding of the drivers behind migration remains weak'' when citation dates to 1993.
\end{reviewbox}

\textbf{Response:} Reframed as a long-recognised, ongoing challenge, with more recent literature acknowledged.

\vspace{0.3cm}
\begin{updatebox}
The theoretical understanding of the drivers behind migration has long been recognised as limited, despite their
significant impact across many countries \citep{227-massey1993theories}. More recent literature continues to highlight
challenges in capturing the complexity and interaction of migration drivers. Historical models of migration have also been
criticised for sample selection bias and insufficient attention to `pull' factors (drivers in potential destination
countries that attract forced migrants) \citep{145-davenport2003leave:}.
\end{updatebox}

\subsubsection*{Section 1.7}

\begin{reviewbox}
First sentence \ldots{} no citations to back this up. Second sentence ``These elements are essential\ldots'' -- says who? for what? Later \ldots{} ``current models struggle'' -- what models?
\end{reviewbox}

\textbf{Response:} The paragraph now distinguishes individual-level / aggregate-flow approaches from
interaction-based ones, justifies the importance of group dynamics and communication, and acknowledges that
group-based movement is not universal.

\vspace{0.3cm}
\begin{updatebox}
Existing models often overlook the role of group dynamics and communication among migrant populations, with many migration
models focusing primarily on individual-level decision-making or aggregate flows rather than interactions between agents.
These elements are essential for understanding certain migration patterns, particularly when people travel together and share
information to guide decisions about potential destinations, as highlighted in studies of migration networks and information
flows. Without accounting for such interactions, these types of models struggle to capture the realities of these forms of
forced migration, leaving a gap in the literature. At the same time, it is important to recognise that group-based movement
is not universal. It is most commonly linked to acute-onset events, such as sudden outbreaks of violence or climate-related
disasters. In other contexts, especially in some waves of displacement, solo migration is more prevalent. For instance,
during the early stages of the 2015 refugee crisis in Lesvos, Greece, the first arrivals were often single men traveling
alone, while family groups became more prominent in later waves as demographics shifted
\citep{unhcr_sea_arrivals_2015,crawley2016unpacking}.
\end{updatebox}

\subsubsection*{Section 1.9}

\begin{reviewbox}
``\ldots{}leads to improved predictive outcomes'' of what? over what? in what specific dimensions of forced migration movement?
\end{reviewbox}

\textbf{Response:} Specified the dimensions (camp population distributions, destination selection, temporal
dynamics), the mechanisms (grouping and communication), and the baseline (Flee). This edit is \emph{merged} with
Reviewer~1's ``social networks'' change into the single §1.8 paragraph shown under Reviewer~1, §1.9.

\vspace{0.3cm}
\begin{updatebox}
Using the case study of Mali in 2012, the research seeks to demonstrate that a more granular understanding of the forced
migrant experience leads to improved predictive outcomes, particularly in modelling camp population distributions,
destination selection, and the temporal dynamics of displacement, through the incorporation of grouping mechanisms and
communication processes, compared to baseline models such as Flee.
\end{updatebox}

\subsubsection*{Section 1.10}

\begin{reviewbox}
Typo in the word ``explain''.
\end{reviewbox}

\textbf{Response:} Corrected (also flagged by Reviewer~1; fixed once).

\subsection*{Literature Review}

\subsubsection*{General Comment}

\begin{reviewbox}
The literature review still feels very weak to set up the importance of the model.
\end{reviewbox}

\textbf{Response:} The literature review has been strengthened. A new ``Motivation'' subsection now establishes
why forecasting forced migration matters and why ABMs are the appropriate methodology, and a new ``Limitations of
existing models'' paragraph consolidates the research gap before Flee and Metafora are introduced. The material
was drawn from the fuller original literature review and condensed to add substance without losing concision.
The added paragraphs (now §2.1--§2.2 and §2.5) are:

\vspace{0.3cm}
\begin{updatebox}
\textbf{§2.1 (Motivation):} Forecasting forced migration is essential for estimating displaced populations, identifying
likely destinations, and supporting humanitarian planning and resource allocation (Edwards 2008; Frydenlund 2021;
Suleimenova et al.\ 2022). Traditional analytical and econometric migration models have contributed substantially to
understanding long-term migration processes but often struggle to represent sudden conflict-induced displacement and
heterogeneous individual decision-making (Bijak et al.\ 2019; de Haas 2021). Consequently, Agent-Based Models (ABMs) have
become an increasingly important methodology within forced migration research because they explicitly represent individual
agents, their interactions, and the emergence of system-level migration patterns (Borshchev \& Filippov 2004; Macal \&
North 2010). These characteristics make ABMs well suited to evaluating humanitarian scenarios and policy interventions under
rapidly evolving crises (Frydenlund 2021).

\medbreak
\textbf{§2.2:} Among forced migration ABMs, the Flee framework has become one of the most influential models, providing a
flexible network-based simulation capable of forecasting refugee movements using publicly available humanitarian datasets
(Groen 2016; Suleimenova et al.\ 2017a, 2022).

\medbreak
\textbf{§2.5 (Limitations of existing models):} Although recent ABMs have substantially advanced forced migration
forecasting, several limitations remain. Existing models often simplify behavioural processes, rely heavily on registered
refugee data, or place limited emphasis on interactions between migrants. These limitations motivate the incorporation of
empirically informed grouping and communication mechanisms within GRASP (Cheong 2023; Suleimenova et al.\ 2017a;
Frydenlund 2021).
\end{updatebox}

\vspace{0.2cm}
\begin{authorbox}
\textbf{Structural note:} the two ``Potential improvements'' paragraphs from v1 (group formation; communication) are
retained as §2.6--§2.7 and now sit under the reorganised headings. No content was cut from them.
\end{authorbox}

\subsubsection*{Section 2.4}

\begin{reviewbox}
2.4 is the most compelling literature review paragraph in the first section. It best establishes the role of this particular model and is situated in more recent literature.
\end{reviewbox}

\textbf{Response:} Thank you. We used §2.4 (Metafora) as the model for tightening and situating the surrounding
literature-review paragraphs.

\subsection*{Results}

\subsubsection*{Section 5.1}

\begin{reviewbox}
``While some deviations will exist between GRASP predictions and empirical camp populations, this does not undermine the value of the model'' -- rather than stating this, can you provide a statistical summary that shows that variations are 1. on par or better than FLEE or other models? and/or 2. statistically insignificant? \ldots{} Similar comment in paragraph 5.3.
\end{reviewbox}

\textbf{Response (partial):} We now report classification performance with pooled and daily precision metrics
(Table~3) computed only at empirical observation dates, which gives a firmer quantitative reference point than the
previous prose. Relative-difference benchmarking against Flee per camp is given in Tables~4--5 and Figs~10--11.

\vspace{0.3cm}
\begin{flagbox}
\textbf{OPEN --- statistical benchmark still needed.} The reviewer specifically asks for a statistical summary
establishing whether GRASP's deviations are (i)~on par with / better than Flee, and/or (ii)~statistically
insignificant against a V\&V benchmark. This is not yet added. Decision needed: report a paired error statistic
(e.g.\ mean/median relative error with confidence intervals, or a paired test of GRASP vs Flee per camp), or
state explicitly why a significance test is not appropriate here. Once chosen, a short updatebox can replace both
the §5.1 and §5.3 defensive sentences.
\end{flagbox}

\subsubsection*{Section 5.6}

\begin{reviewbox}
About the limitations re: utility function reliance -- what impact does this have? The explanation about lengthy run times falls a bit flat rather than focusing on what we can know from this exercise.
\end{reviewbox}

\textbf{Response:} Refocused from run-time onto the modelling implication: how parameter sensitivity in the utility
function affects destination/camp-allocation predictions, and why this motivates qualitative calibration of the
micro-level behavioural assumptions.

\vspace{0.3cm}
\begin{updatebox}
The main limitation of the GRASP model lies in its reliance on a calibrated utility function for destination selection.
While parameter optimisation was performed at higher \textit{frac} values due to computational constraints, this introduces
uncertainty when the model is scaled to full-resolution simulations, as parameter sensitivity may vary with population size.
Consequently, the utility function primarily affects the allocation of refugees between camps rather than the total number
of refugees predicted. The utility-based approach nevertheless consistently outperformed the simpler alternative of selecting
the nearest available destination, supporting its inclusion despite the remaining inaccuracies in camp allocation
(Fig.~\ref{emp0}). It is therefore important not to rely solely on aggregate model output for validation, as \citet{He-In}
note that this approach is insufficient. Instead, evaluating the underlying micro-processes is essential, motivating the
incorporation of qualitative data from forced migrants to better inform the utility-based decision mechanism.
\end{updatebox}

% ======================================================
% ============ AUTHOR-INITIATED CHANGES ================
% ======================================================
\newpage
\section*{Author-Initiated Changes (Not Requested by Reviewers)}
\label{sec:author}
\addcontentsline{toc}{section}{Author-Initiated Changes (Not Requested by Reviewers)}

The following changes between the previous submission (v1) and this version were made at the authors'
initiative rather than in response to a specific comment. Each is listed with its rationale. Items~A1--A2
account for essentially all of the numerical differences (abstract percentages, Tables~3--5, and the inline
accuracy/precision values).

\subsection*{A. Reproducibility and computational engineering}

\begin{authorbox}
\textbf{A1 --- Fixed-seed re-run of all simulations.} All results were regenerated with a fixed random seed to
remove run-to-run stochastic variation, so every table and figure is now deterministic and reproducible.
\emph{Rationale:} reproducibility, and a direct response to the Fig.~6 replication concern. \emph{Effect:}
headline camp precision 75.1\% $\rightarrow$ 77.85\%; grouping precision gain 46\% $\rightarrow$ 48.8\%;
communication relative-difference reduction 56\% $\rightarrow$ 61.26\%; total-refugee accuracy 98.4\%
$\rightarrow$ 98.99\%; IDP classification 85.5\% $\rightarrow$ 87.95\%; and all values in Tables~3, 4 and 5.
\end{authorbox}

\begin{authorbox}
\textbf{A2 --- Metrics computed only at empirical observation dates.} All accuracy, precision and
relative-difference metrics are now calculated only on empirical UNHCR observation dates; interpolated points are
shown in figures (marked by vertical dash-dotted lines) but never scored. New/added methodological text:
§3.12 (summary precision defined as the arithmetic mean of daily precision; pooled precision defined), §3.14
(metrics at empirical dates only), §3.15 (dash-dotted observation markers), and §5.8 (why pooled metrics differ
slightly from daily). \emph{Rationale:} interpolated points artificially smooth and inflate metrics; scoring only
on measured dates is more honest. \emph{Effect:} contributes to the number changes in A1.
\end{authorbox}

\begin{authorbox}
\textbf{A3 --- Windows compatibility fix.} A path-handling bug was fixed so the model now runs on Windows.
\emph{Rationale:} broaden usability and reproducibility across operating systems.
\end{authorbox}

\begin{authorbox}
\textbf{A4 --- Dependency updates.} Updated to current versions of key dependencies (e.g.\ \texttt{pandas}).
\emph{Rationale:} maintainability, security, and reproducibility on current environments.
\end{authorbox}

\begin{authorbox}
\textbf{A5 --- Checkpointing / crash recovery.} Added alternating-slot checkpointing (atomic \texttt{.tmp} +
\texttt{os.replace}) so long runs can be resumed after an interruption without corruption. \emph{Rationale:}
robustness for long simulations; supports reproducibility of the full run.
\end{authorbox}

\begin{authorbox}
\textbf{A6 --- Simpler feature toggling.} The grouping, communication and roulette modules can now be switched on
or off more easily. \emph{Rationale:} cleaner, reproducible ablation studies (the ``without grouping'' / ``without
communication'' variants in Tables~4--5 and Figs~12--13).
\end{authorbox}

\subsection*{B. Model / calibration changes}

\begin{authorbox}
\textbf{B1 --- Bias term added to the utility function.} A bias term $\beta_0$ was added to the utility equation
(now Eq.~12), with $\beta_0 = 150$ in Table~2, and $\beta_{F,i}$ was changed from $1/2$ to $1$. \emph{Rationale:}
$\beta_0$ acts as a parameter of ``irrationalism'' that raises the likelihood of rarer choices; it also connects
naturally to Reviewer~1's §4.20 point about learnable $\beta$ parameters. \emph{Effect:} recalibration consistent
with the re-run.
\end{authorbox}

\subsection*{C. Presentation and structure}

\begin{authorbox}
\textbf{C1 --- Example camps switched to Tabareybarey/Mbera.} The empirical-comparison figures and text (Fig.~5,
Fig.~10, Fig.~11; §5.3, §5.4, §5.14--§5.16) now use Tabareybarey and Mbera instead of Mentao and Abala.
\emph{Rationale:} under the re-run these two better illustrate the contrasting over-prediction (transit-hub) and
under-prediction (large-camp convergence) behaviours.
\end{authorbox}

\begin{authorbox}
\textbf{C2 --- 7-day moving-average plotting.} Simulated series are now shown as a 7-day moving average, with a
note distinguishing smoothed from unsmoothed outputs. \emph{Rationale:} clearer visualisation of transition
dynamics without hiding the underlying spikes.
\end{authorbox}

\begin{authorbox}
\textbf{C3 --- Removed v1 §1.8 (GRASP ``innovative features'' paragraph).} The intro paragraph that listed GRASP's
innovations (advanced grouping/communication, IDP/returnee integration, empirical foundation for stochastic
variation) was removed. \emph{Rationale:} redundant with the expanded literature review and methodology; tightens
the introduction.
\end{authorbox}

\begin{authorbox}
\textbf{C4 --- Removed a sentence in §1.4} (``These efforts focus on bridging the gap between statistical
information and the understanding of causal systems''). \emph{Rationale:} tightening; not needed after the
granularity rephrase.
\end{authorbox}

\begin{authorbox}
\textbf{C5 --- Figure and equation renumbering.} Removing the confusion-matrix figure (Reviewer~1) shifted all
later figure numbers down by one; inserting the wealth equation as Eq.~(2) shifted all later equation numbers up
by one. \emph{Rationale:} mechanical consequence of the two accepted edits --- flagged so cross-references are
checked on recompile.
\end{authorbox}

\begin{authorbox}
\textbf{C6 --- Added §5.20 (feature importance is location-specific).} A new paragraph notes that removing
communication most affects Bobo, removing grouping most affects Goudoubo/Mentao/Ouagadougou, and the roulette
step most affects Mbera. \emph{Rationale:} supports and interprets the ablation results.
\end{authorbox}

\begin{authorbox}
\textbf{C7 --- Minor Discussion rewording} in §6.9, §6.11 and §6.12 (e.g.\ ``over 80\% accuracy'' $\rightarrow$
``over 70\% average precision'' at §6.8, and an added explanation of the temporary precision dip and end-of-run
decline). \emph{Rationale:} accuracy with the re-run figures and clearer phrasing.
\end{authorbox}

\begin{authorbox}
\textbf{C8 --- Appendices.} Appendix figures were regenerated to match the reseeded run. No other appendix changes
are documented, per author instruction.
\end{authorbox}

% ======================================================
% ================= OPEN ITEMS =========================
% ======================================================
\newpage
\section*{Outstanding Items Requiring Author Input}
\label{sec:open}
\addcontentsline{toc}{section}{Outstanding Items Requiring Author Input}

The red boxes above are collected here for convenience. These need a decision or a value from you before final
submission.

\begin{flagbox}
\textbf{Reviewer responses still empty / partial:}
\begin{itemize}
  \item \textbf{Fig.~6 (repository / replication)} --- confirm the missing CSV(s) are committed, do a clean-clone
        reproduction test, and state the random seed in the README.
  \item \textbf{§6.12 (computational cost)} --- insert the actual runtime and hardware for a full run.
  \item \textbf{Reviewer~2 §5.1 (statistical summary)} --- add a statistical benchmark vs Flee, or justify why a
        significance test is not appropriate.
  \item \textbf{§3.13/§3.14} --- decide whether to add a one-sentence justification of precision-over-accuracy
        (currently fully removed).
  \item \textbf{Fig.~11 (old Fig.~12) lower subplot} --- cap $y$ at 1.25 as requested, then confirm.
\end{itemize}
\end{flagbox}

\begin{flagbox}
\textbf{Leftover errors to fix in the manuscript (found during the diff):}
\begin{itemize}
  \item §4.9 footnote is still the placeholder ``Ref''.
  \item §5.5 references ``Section 66.12'' (should be 6.12, ideally a link).
  \item §6.10 contains an unresolved ``Section 5 ??'' cross-reference.
  \item §6.12 contains the typo ``rovide'' (should be ``provide'').
  \item The sentence ``It rarely overestimates \ldots{} conservative estimator of camped refugee numbers'' is
        duplicated across §6.8 and §6.9.
  \item The Results numbering resets: §5.20 is followed by two paragraphs numbered §5.1 and §5.2 (should be §5.21,
        §5.22).
  \item §5.6 contains the typo ``t is therefore'' (should be ``It is therefore'') --- corrected in the updatebox
        above.
  \item Verify the ``Network (2017)'' bibliography entry now renders as the DNV~GL report after recompilation.
  \item Check all figure/equation \texttt{\textbackslash ref}s after the renumbering in item~C5.
\end{itemize}
\end{flagbox}

\end{document}
